% Segment Tree Beats - range chmin + range add + range sum

Supports range add, range chmin ($a_i = \min(a_i, x)$), range sum, range max. $O((N+Q)\log N)$ amortized. For range chmax, mirror the struct by swapping max/min roles (maintain \texttt{mn1, mn2, cnt\_mn} and \texttt{apply\_chmax}).

\begin{lstlisting}[language=C++]
struct SegTreeBeats {
    static constexpr int MAXN = 200005;
    static constexpr long long INF = (1LL << 60);
    int n;
    long long sum[MAXN << 2], mx1[MAXN << 2], mx2[MAXN << 2], add[MAXN << 2];
    int cnt[MAXN << 2], len[MAXN << 2];

    void pull(int p) {
        sum[p] = sum[p << 1] + sum[p << 1 | 1];
        if (mx1[p << 1] == mx1[p << 1 | 1]) {
            mx1[p] = mx1[p << 1];
            cnt[p] = cnt[p << 1] + cnt[p << 1 | 1];
            mx2[p] = max(mx2[p << 1], mx2[p << 1 | 1]);
        } else if (mx1[p << 1] > mx1[p << 1 | 1]) {
            mx1[p] = mx1[p << 1];
            cnt[p] = cnt[p << 1];
            mx2[p] = max(mx2[p << 1], mx1[p << 1 | 1]);
        } else {
            mx1[p] = mx1[p << 1 | 1];
            cnt[p] = cnt[p << 1 | 1];
            mx2[p] = max(mx1[p << 1], mx2[p << 1 | 1]);
        }
    }

    void apply_add(int p, long long v) {
        sum[p] += v * len[p];
        mx1[p] += v;
        if (mx2[p] != -INF) mx2[p] += v;
        add[p] += v;
    }

    void apply_chmin(int p, long long x) {
        // precondition: mx2[p] < x < mx1[p]
        sum[p] -= (mx1[p] - x) * cnt[p];
        mx1[p] = x;
    }

    void push(int p) {
        if (add[p] != 0) {
            apply_add(p << 1, add[p]);
            apply_add(p << 1 | 1, add[p]);
            add[p] = 0;
        }
        if (mx1[p << 1] > mx1[p]) apply_chmin(p << 1, mx1[p]);
        if (mx1[p << 1 | 1] > mx1[p]) apply_chmin(p << 1 | 1, mx1[p]);
    }

    void build(int p, int l, int r, vector<long long>& a) {
        len[p] = r - l + 1;
        add[p] = 0;
        if (l == r) {
            sum[p] = mx1[p] = a[l];
            mx2[p] = -INF;
            cnt[p] = 1;
            return;
        }
        int m = (l + r) >> 1;
        build(p << 1, l, m, a);
        build(p << 1 | 1, m + 1, r, a);
        pull(p);
    }

    void range_add(int p, int l, int r, int ql, int qr, long long v) {
        if (qr < l || r < ql) return;
        if (ql <= l && r <= qr) { apply_add(p, v); return; }
        push(p);
        int m = (l + r) >> 1;
        range_add(p << 1, l, m, ql, qr, v);
        range_add(p << 1 | 1, m + 1, r, ql, qr, v);
        pull(p);
    }

    void range_chmin(int p, int l, int r, int ql, int qr, long long x) {
        if (qr < l || r < ql || mx1[p] <= x) return;
        if (ql <= l && r <= qr && mx2[p] < x) { apply_chmin(p, x); return; }
        push(p);
        int m = (l + r) >> 1;
        range_chmin(p << 1, l, m, ql, qr, x);
        range_chmin(p << 1 | 1, m + 1, r, ql, qr, x);
        pull(p);
    }

    long long range_sum(int p, int l, int r, int ql, int qr) {
        if (qr < l || r < ql) return 0;
        if (ql <= l && r <= qr) return sum[p];
        push(p);
        int m = (l + r) >> 1;
        return range_sum(p << 1, l, m, ql, qr) + range_sum(p << 1 | 1, m + 1, r, ql, qr);
    }

    long long range_max(int p, int l, int r, int ql, int qr) {
        if (qr < l || r < ql) return -INF;
        if (ql <= l && r <= qr) return mx1[p];
        push(p);
        int m = (l + r) >> 1;
        return max(range_max(p << 1, l, m, ql, qr), range_max(p << 1 | 1, m + 1, r, ql, qr));
    }

    // public wrappers (1-indexed array a[1..n])
    void init(vector<long long>& a) {
        n = (int)a.size() - 1;
        build(1, 1, n, a);
    }
    void add(int l, int r, long long v) { range_add(1, 1, n, l, r, v); }
    void chmin(int l, int r, long long x) { range_chmin(1, 1, n, l, r, x); }
    long long query_sum(int l, int r) { return range_sum(1, 1, n, l, r); }
    long long query_max(int l, int r) { return range_max(1, 1, n, l, r); }
};
\end{lstlisting}